\section{Warm-Start Role}
\label{sec:warmstart}

\subsection{Why It Might Help}

MCMC and ReCom-style workflows need an initial plan. A \metis{} warm-start is
standard because it is fast and tends to have low edge cut. A \bfsg{} warm-start
is different: it is also fast, but its boundaries are produced by frontier
growth rather than multilevel cut refinement. That can be useful when the goal
is to test sensitivity to starting geography.

\subsection{What Is Not Yet Proven}

The paper should not yet claim that \bfsg{} improves burn-in, effective sample
size, or autocorrelation decay. Those claims require paired chains from METIS
and BFS Growth starts, identical proposal kernels, identical seeds where
applicable, archived traces, and diagnostic summaries.

\subsection{Warm-Start Test Plan}

A proper warm-start package should record:
\begin{itemize}
\item initial plan EC, PP, population deviation, and district contiguity;
\item chain proposal type and all proposal seeds;
\item accepted/rejected proposal ledger;
\item EC and compactness traces;
\item effective sample size and R-hat when multiple chains are run; and
\item final comparison between METIS-started and BFS-started chains.
\end{itemize}

Until that package exists, the practical recommendation is conservative:
\metis{} remains the default warm-start, and \bfsg{} is an alternate starting
condition for robustness checks.

\subsection{CLI Shape}

\begin{verbatim}
bisect state --state NC --year 2020 \
  --structure bfs-growth --search single
\end{verbatim}

The resulting plan can be used wherever the platform accepts a previously
generated plan as an initial condition. The important audit fields are the base
seed, node path, BFS node seed, BFS split seed, adjacency hash, and population
hash.
